\slajdid{02-s013}
\begin{frame}[label=02-s013]
\gusttekst
\begin{columns}[c,onlytextwidth]
\begin{column}{.74\textwidth}
Na osnovu osobina logičkih funkcija sada je iz funkcionalne tabele moguće napisati logičke funkcije. Posmatrajući kada funkcija ima vrednost logičke jedinice iskazi rečima osobina su redom od vrha funkcionalne table:
\begin{enumerate}
\item Funkcija ima vrednost logičke jedinice ako promenjiva C ima vrednost logičke nule I promenjiva B vrednost logičke jedinice I promenjiva A vrednost logičke jedinice
\item ILI - funkcija ima vrednost logičke jedinice ako promenjiva C ima vrednost logičke jedinice I promenjiva B vrednost logičke nule I promenjiva A vrednost logičke jedinice
\item ILI - itd\ldots
\end{enumerate}
Uočite da su iskazi povezani I i ILI logičkim funkcijama i da kada je potrebno da promenjiva ima vrednost logičke nule njena komplementna vrednost ima vrednost logičke jedinice. Prema tome možemo iz funkcionalne tabele napisati izraz
\[
F=\bar CBA+C\bar BA+CB\bar A
\]
\end{column}
\begin{column}{.23\textwidth}
\begingroup
\setlength{\tabcolsep}{3pt}
\renewcommand{\arraystretch}{1.2}
\par\noindent
\begin{tabularx}{\linewidth}{*{4}{>{\centering\arraybackslash}X}}
\toprule C&B&A&F\\\midrule
0&0&0&0\\
0&0&1&0\\
0&1&0&0\\
\textcolor{DEcrvena}{0}&\textcolor{DEcrvena}{1}&\textcolor{DEcrvena}{1}&\textcolor{DEcrvena}{1}\\
1&0&0&0\\
\textcolor{DEcrvena}{1}&\textcolor{DEcrvena}{0}&\textcolor{DEcrvena}{1}&\textcolor{DEcrvena}{1}\\
\textcolor{DEcrvena}{1}&\textcolor{DEcrvena}{1}&\textcolor{DEcrvena}{0}&\textcolor{DEcrvena}{1}\\
1&1&1&0\\\bottomrule
\end{tabularx}\par
\endgroup

\end{column}
\end{columns}
\end{frame}
